% Transcription — fonds Alexandre Grothendieck, shelfmark 22, batch 1 (pages 1-20).
% Established from the facsimile archives/batches/22/batch-01.pdf (a mirror of
% https://grothendieck.umontpellier.fr/22.pdf), per the skill
% transcribe-grothendieck. The batch file carries no cover sheet: PDF page N is
% archivists' page N. Checked against the pencilled numbers bottom-left — "1"
% through "8" on the eight offprint leaves, then "9" … "17" on the typescript,
% "18" on the blank cover and "19", "20" on the two manuscript leaves — the
% alignment holds end to end.
%
% Pass: Opus 5 (claude-opus-5), 2026-09-19 — first pass, unchecked against the pages by a human.
%
% What the batch is, and the sorting it required. The inventory announces three
% kinds of paper and all three are here, in three consecutive blocks.
%
%   1-8    Two offprints, printed, not in his hand: J.-P. Jouanolou's two
%          Comptes rendus notes of 1963 — « Étude différentielle des anneaux
%          locaux réguliers : applications » (t. 257, pp. 2948-2950, séance du
%          4 novembre 1963) on pages 1-3 with its blank verso on page 4, and
%          « Sur les modules de différentielles » (t. 256, pp. 842-844, séance
%          du 14 janvier 1963) on pages 5-7 with its blank verso on page 8.
%          Printed matter is not transcribed. Pages 1, 2, 3 and 5 carry
%          annotations in his hand and appear for those annotations alone,
%          each anchored by a \note saying which printed statement it sits
%          beside; the printed text itself is not reproduced. Pages 4, 6, 7
%          and 8 carry nothing in his hand and are simply absent.
%   9-17   The annotated typescript, nine leaves, paginated by him 24 to 32 in
%          ink at the top right (the archivists' 9 to 17 are pencilled bottom
%          left). It is a continuous run of the numéro 3 of a chapter: the
%          theorem 3.6 on formally simple algebras, the definition 3.7 of a
%          Cohen algebra, and the corollaries 3.8 to 3.13. It is addressed to
%          a reader in the second person — « j'espère que tu pourras la rendre
%          plus lisible », « que tu regardes cela » — so the typescript is a
%          draft sent out, and the annotations are his corrections to it.
%   18     An otherwise blank leaf inscribed « Questions » in his hand at the
%          top right: the cover of the run that follows. It gets no \page; its
%          word opens the last section.
%   19-20  Two leaves of manuscript, ink, the open questions 1) to 4), each
%          leaf carrying besides the numbered questions a dense block of
%          later additions written diagonally up the left margin.
%
% The typed section numbers of the typescript are struck and overtyped at
% nearly every occurrence — the same correction made mechanically down the run
% — and what stands under them is not always decidable on the film. They are
% transcribed as 3.N throughout, which is what the untouched occurrences in
% the body of the prose (« grâce à 3.1 », « en vertu de 3.5 », « de 3.9 »)
% and the handwritten corrections both give; the point is flagged once, at
% page 10.
%
% What is NOT here. No attempt was made to identify the chapter of EGA IV the
% typescript became, nor to reconcile its numbering with the printed text: the
% numbering is the draft's own. The printed Jouanolou notes are described,
% never transcribed.
\documentclass[11pt,a4paper]{article}
% Le préambule commun à toutes les transcriptions du fonds Grothendieck.
%
% Il tient en une idée : **l'incertitude se compose, elle ne se résout pas.**
% Une transcription qui lisse un mot douteux a détruit la seule chose pour
% laquelle on la faisait — un lecteur doit pouvoir savoir, à chaque endroit,
% ce qui était sur le feuillet et ce qui a été ajouté. Les six macros qui
% suivent sont donc l'appareil critique, et non de la décoration.
%
% Les mêmes commandes sont reconnues par `scripts/render.mjs`, qui produit la
% vue de lecture du site. Une macro ajoutée ici doit l'être aussi là-bas,
% faute de quoi le rendu échouera bruyamment — ce qui est voulu : un rendu
% silencieusement faux est pire qu'un rendu absent.

\ProvidesPackage{grothendieck}[2026/08/08 Transcriptions du fonds Grothendieck]

\usepackage{amsmath,amssymb,amsthm}
\usepackage{mathrsfs}
\usepackage{stmaryrd}
% Les diagrammes commutatifs sont partout dans ce fonds : c'est la langue
% même de la géométrie algébrique de Grothendieck, et une transcription qui
% les rendrait en prose ne transcrirait rien.
\usepackage{tikz-cd}
% Français par défaut — c'est la langue des feuillets — mais l'anglais est
% chargé aussi : la traduction et le résumé sont des documents anglais, et la
% typographie française (espaces avant les deux-points) y serait une faute.
% Ces documents basculent par \selectlanguage{english} en tête de corps.
\usepackage[english,french]{babel}
\usepackage{fontspec}
\usepackage{geometry}
\geometry{a4paper,margin=25mm,marginparwidth=28mm}
\usepackage{marginnote}
\usepackage{xcolor}
% ulem, not soul. `soul`'s \st and \ul are built on a character-by-character
% scan that XeLaTeX's font machinery breaks: the document fails with an
% undefined control sequence rather than degrading. ulem does the same two jobs
% and is Unicode-engine safe. `normalem` stops it hijacking \emph, which the
% transcriptions use for Grothendieck's own emphasis.
\usepackage[normalem]{ulem}
\usepackage{hyperref}
\hypersetup{colorlinks=true,linkcolor=black,urlcolor=black,pdfborder={0 0 0}}

% --- Métadonnées du lot -----------------------------------------------------
% Reprises telles quelles par le script de rendu, qui les affiche en tête de
% la vue de lecture. Elles fixent ce que le fichier prétend couvrir : sans
% elles, une transcription flotte sans point d'attache dans un fonds de seize
% mille feuillets.

\newcommand{\folder}[1]{\gdef\grothFolder{#1}}
\newcommand{\batch}[1]{\gdef\grothBatch{#1}}
\newcommand{\leaves}[2]{\gdef\grothFirst{#1}\gdef\grothLast{#2}}
\newcommand{\foldertitle}[1]{\gdef\grothFoldertitle{#1}}
\gdef\grothFolder{?}\gdef\grothBatch{?}\gdef\grothFirst{?}\gdef\grothLast{?}\gdef\grothFoldertitle{}

% --- Le résumé --------------------------------------------------------------
%
% Il ouvre la lecture modernisée, sous le titre « Résumé », et il est composé
% en petit. Ce n'est pas un ornement : c'est le seul passage du document qui
% ne soit pas des mathématiques, et un lecteur doit pouvoir le reconnaître
% avant de l'avoir lu — puis le sauter s'il connaît déjà le sujet, ou s'y
% arrêter s'il ne le connaît pas. Le filet qui le suit marque la reprise du
% corps.

\newenvironment{resume}
  {\small\setlength{\parskip}{0.45em}}
  {\par\medskip\hrule\medskip}

% --- Mots-clés ---------------------------------------------------------------
%
% En anglais, en clôture du résumé de la lecture modernisée : le vocabulaire
% moderne sous lequel chercher ce que les feuillets construisent. Le manifeste
% du site (`npm run manifest`) les extrait pour étiqueter la cote entière —
% cette ligne est la source unique des tags, il n'y a pas de fichier à côté.
\newcommand{\keywords}[1]{\par\smallskip\noindent
  {\footnotesize\textsc{Keywords} --- \textit{#1}}\par}

% --- L'appareil critique ----------------------------------------------------

% Le numéro de feuillet, en marge. C'est l'ancre qui permet de régler un
% désaccord sur une formule en regardant *un* feuillet plutôt qu'une chemise
% de six cents. Le site s'en sert aussi pour tourner les pages du fac-similé
% quand on fait défiler la transcription.
\newcommand{\leaf}[1]{%
  \marginnote{\footnotesize\color{gray!70}\textsf{#1}}%
  \label{leaf:#1}\ignorespaces}

% Illisible : on ne devine pas. Un mot inventé qui se lit comme les autres est
% le pire résultat possible.
\newcommand{\ill}{\textcolor{red!60!black}{[\,\ldots\,]}}

% Lecture douteuse : le mot est donné, et signalé comme incertain.
\newcommand{\uncertain}[1]{\uline{#1}}

% Ajout de l'éditeur — une lettre manquante, un mot sous-entendu.
\newcommand{\add}[1]{\textcolor{blue!50!black}{[#1]}}

% Biffé par Grothendieck lui-même. Ce qu'il a rayé fait partie du document :
% une rature dit souvent où la pensée a changé de direction.
\newcommand{\struck}[1]{\sout{#1}}

% Note du transcripteur, sur l'état matériel du feuillet.
\newcommand{\note}[1]{\par\smallskip{\small\color{gray!60!black}\textit{[#1]}}\par\smallskip}

% Note marginale de Grothendieck, distincte du corps du texte.
\newcommand{\marginal}[1]{\par\smallskip\noindent\hspace{1em}%
  {\small\color{gray!60!black}#1}\par\smallskip}

% --- Titre ------------------------------------------------------------------

\AtBeginDocument{%
  \begin{center}
    {\large\bfseries Fonds Alexandre Grothendieck --- cote n\textsuperscript{o}~\grothFolder}\\[2pt]
    {\itshape\grothFoldertitle}\\[4pt]
    {\small feuillets \grothFirst--\grothLast\ (lot \grothBatch)}\\[2pt]
    {\footnotesize\color{gray!60!black}Transcription établie d'après le fac-similé numérisé
     par l'Université de Montpellier.\\
     Les lectures incertaines sont soulignées, les illisibles notées [\,\ldots\,],
     les ajouts entre crochets.}
  \end{center}
  \vspace{1em}}

\endinput


\folder{22}
\batch{1}
\pages{1}{20}
\dating{1963-[vers 1967]}
\watermark{Édition de démonstration}
\foldertitle{[EGA IV]. [Chapitre] 0 IV. Cohen, différentielles, algèbres formellement simples (Compléments) : tirés à part (1963), tapuscrit annoté (s.d.), notes manuscrites (s.d.).}

\begin{document}

\section{Tirés à part annotés : les deux notes de Jouanolou (1963)}

\page{1}
\note{Tiré à part, imprimé : la note de J.-P. Jouanolou « Étude différentielle
des anneaux locaux réguliers : applications », Comptes rendus, t. 257,
p. 2948-2950, séance du 4 novembre 1963. Le texte imprimé n'est pas de sa main
et n'est pas transcrit ; ne sont relevées ici que ses annotations, avec
l'indication de l'énoncé qu'elles visent.}

\marginal{\uncertain{appel}}
\note{en marge de gauche, à hauteur de la première phrase du nº 1, celle qui
rappelle les critères différentiels de régularité de Nakai, Suzuki et Mount.}

\marginal{\ill{} \uncertain{hypothèse}}
\note{en marge de droite, à la même hauteur. Dans la ligne imprimée il biffe le
mot « condition » de « une condition de réduction analytique » et porte deux
signes d'insertion ; l'annotation marginale est vraisemblablement le mot de
remplacement, mais le premier trait ne se lit pas.}

\marginal{$A/\mathfrak{a}$}
\note{en marge de gauche, en regard de la Proposition 1, dont l'hypothèse porte
sur $D(A)$ plat et $\mathfrak{a}$ un idéal tel que $A/\mathfrak{a}$ soit réduit ;
il entoure au crayon le mot « réduit » de l'énoncé, et une accolade réunit la
Proposition 1 et le point (d), $D(A/\mathfrak{a})$ libre.}

\marginal{\ill{}}
\note{trois lignes courtes en marge de droite, en regard du Corollaire 2
[Kunz] — celui qui caractérise la régularité d'une localité réduite par la
liberté du module de différentielles absolu $D(A)$. Le premier mot pourrait se
lire « contenu », la suite ne se laisse pas lire.}

\marginal{\uncertain{assez} rare !}
\note{en marge de droite, en regard du Corollaire 3 [Suzuki, Mount], dont
l'hypothèse est que le séparé $D'(A)$ soit de type fini.}

\marginal{séries formelles ??}
\note{en marge de droite, en regard de la Proposition 2, sur le lieu singulier
d'une $A$-algèbre affine ; une accolade réunit les Propositions 2 et 3. La
question porte sur le passage des séries convergentes du Corollaire 1 aux
séries formelles.}

\marginal{\ill{}}
\note{un mot en marge de droite, en regard de la Proposition 3, sur le lieu
singulier fermé d'un schéma de type fini sur un schéma en groupes régulier ;
il se termine par « -ique » et n'a pas pu être lu.}

\page{2}
\note{Suite du même tiré à part, page 2 de la note imprimée.}

\marginal{connu}
\note{en marge de droite, au bout d'un long trait vertical qui couvre tout le
bloc allant du Lemme sur le noyau de $\hat{A} \otimes_A D(A) \to D(\hat{A})$
jusqu'à la Proposition 6, c'est-à-dire les énoncés de séparabilité de
$K_{\hat{A}}$ sur $K_A$ et le passage par le théorème de structure de Cohen.}

\note{Il entoure à l'encre la condition $(h_2)$ du Théorème, « $D(A)$ est
libre ».}

\marginal{faux \ill{} $A \to B$ \ill{} ??}
\note{en marge de droite, en regard du Lemme du nº 3 — celui qui, de $A$ local
nœthérien analytiquement intègre et $K_{\hat{A}}$ séparable sur $K_A$, conclut
que $B = C_{\mathfrak{p}}$ est analytiquement réduit. L'annotation est écrite en
oblique sur quatre lignes brèves et n'est lisible qu'en partie : le premier mot
est « faux », il y figure une flèche $A \to B$, et elle s'achève sur deux points
d'interrogation.}

\page{3}
\note{Suite du même tiré à part, page 3 de la note imprimée.}

\marginal{$C \supset A$ !!!}
\note{en grand, en haut de la page, au bout d'un trait qui descend vers la
première ligne imprimée ; il entoure le $C$ de « toute $A$-algèbre intègre qui
est localisée d'une $A$-algèbre finie $C$ ». L'objection porte sur ce que
l'énoncé suppose tacitement $A$ contenu dans $C$.}

\marginal{faux}
\note{en marge de droite, au bout d'un long trait horizontal qui souligne la
parenthèse « (Pour appliquer le lemme, on remarque que $h_1$ et $h_2$ sont
stables par localisation.) ».}

\marginal{connu (\uncertain{Nagata})}
\note{en marge de droite, en regard de la Proposition 7 : finitude de la
fermeture intégrale de $A$ dans toute extension finie $L$ de $K_A$.}

\marginal{?}
\note{en marge de droite, en regard de la phrase qui applique la Proposition 7
aux localités, aux localités analytiques et aux algèbres de Hopf régulières sur
un corps parfait ; il souligne « algèbres de Hopf régulières sur un corps
parfait ».}

\marginal{connu}
\note{en marge de droite, en regard du dernier alinéa, celui qui tire du lemme 2
de Zariski-Samuel qu'une localité analytique normale vérifiant le Corollaire 2
est analytiquement normale.}

\page{5}
\note{Second tiré à part, imprimé : la note de J.-P. Jouanolou « Sur les modules
de différentielles », Comptes rendus, t. 256, p. 842-844, séance du 14 janvier
1963. Même règle que ci-dessus : seules les annotations sont transcrites.}

\marginal{non vérifié}
\note{en marge de gauche, souligné, en regard de la définition du nº 1 : $B$ est
dite quasi-séparable lorsque $D_A(B) = 0$. Le trait vise le caractère local de
la propriété, affirmé dans la phrase suivante.}

\marginal{?}
\note{en bas à gauche, en regard de la dernière ligne, « ce qui exprime que $B$
n'est pas ramifié sur $A$ », qui traduit la relation $\mathfrak{n} = \mathfrak{m}B$.}

\section{Tapuscrit annoté : le nº 3, algèbres de Cohen et algèbres formellement simples}

\note{Les neuf feuillets qui suivent forment une suite continue, paginée par lui
24 à 32 à l'encre en haut à droite. Le corps dactylographié est transcrit ;
les surfrappes de la machine deviennent \struck{}, les ajouts manuscrits
interlinéaires \add{}, les notes qu'il porte en marge \marginal{}. Le chiffre
qui ouvre chaque numéro d'énoncé est frappé puis repassé à l'encre, et ce qui
est dessous ne se décide pas sur le film ; les renvois dactylographiés du
corps du texte (« grâce à 3.1 », « en vertu de 3.5 », « dans 3.9 ») donnent
3.N, et c'est cette forme qui est retenue partout ici.}

\page{9}\note{sa page 24}

\marginal{\ill{} inutile, résulte du fait \struck{\ill{}} que $IB$ (ou $IC$)
\uncertain{top.} nilpotent.}

\ldots\ contenant $IB$ resp. \uncertain{$IC$}, enfin on suppose $IB$ et $IC$
fermés (condition automatiquement vérifiée si $B$, $C$ sont nœthériens).
Soient $B_0 = B/IB$, $C_0 = C/IC$, $u$ : \add{$B \to C$} un homomorphisme
continu de $A$-algèbres, $u_i : B_i \to C_i$ l'homomorphisme
\struck{induit} déduit par passage au quotient (j'ai oublié d'introduire un
système fondamental d'idéaux ouverts $I_i \subset I_0$ \add{$= I$} dans $A$,
et les $B_i$ et $C_i$ en conséquence \add{enfin ce sont tous les $I_iB$ et
$I_iC$ qui doivent être fermés}). Conditions équivalentes :
\note{la phrase « enfin ce sont tous les $I_iB$ et $I_iC$ qui doivent être
fermés » est frappée en interligne et rattachée à la parenthèse par un trait
de sa main ; il souligne « contenant $IB$ resp. $IC$ » et corrige au-dessus le
$T$ d'un « TC » frappé pour « IC ».}

\begin{enumerate}
  \item[(i)] $u$ est \struck{un isomorphisme} bijectif, et $u_0$ est un
    homéomorphisme.
  \item[(i bis)] $u$ est un isomorphisme d'algèbres topologiques
  \item[(ii)] Pour tout $i$, $u_i : B_i \to C_i$ est bijectif, $u_0$ est un
    homéomorphisme.
  \item[(ii bis)] pour tout $i$, munissant $B_i$ et $C_i$ des filtrations
    $I$-adiques, l'homomorphisme $\mathrm{gr}_I(B_i) \to \mathrm{gr}_I(C_i)$
    induit par $u$ est bijectif, $u_0$ un homéomorphisme.
\end{enumerate}
\note{les étiquettes (i) et (i bis) portent l'une et l'autre une surcharge,
un « bis » biffé sur la première ; un long trait courbe monte de la seconde
vers la première. Il entoure à l'encre les quatre occurrences de « $u_0$ est
un homéomorphisme » et porte devant chacune une flèche montante : la clause
est à détacher de la liste.}

\marginal{\ill{}}
\note{une note brève de sa main, écrite verticalement dans la marge de gauche
à hauteur de (ii bis) ; elle n'a pas pu être lue.}

Enfin, lorsque les \uncertain{$C_i$} sont plats sur les $A_i$, ces conditions
équivalent également à

\begin{enumerate}
  \item[(iii)] $u_0 : B_0 \to C_0$ est un \struck{isomorphisme}
    homéomorphisme.
\end{enumerate}

L'équivalence de (i) et (i bis) résulte trivialement du fait que les topologies
de $B$, $C$ sont \struck{adiques} définies par les puissances d'idéaux
contenant $IB$ resp. $IC$. \struck{\ill{}} Évidemment (i bis) $\Rightarrow$
(ii), pour (ii) $\Rightarrow$ (i bis) on note que les applications canoniques
$B \to \varprojlim B_i$ et $C \to \varprojlim C_i$ sont bijectives grâce au
fait que les $I_iB$ et $I_iC$ sont fermés (cf rectification à $0_{7.2.4}$).
L'équivalence de (ii) et (ii bis), \struck{qu} mise pour mémoire, provient du
fait que ($I = I_0$ étant un idéal de définition, donc nilpotent dans les
$A_i$) les filtrations mises sur $B$, $C$ sont finies. Enfin, reste à prouver
que lorsque \struck{\ill{}} $C_i$ est plat sur $A_i$, alors $u_0$ bijectif
implique $u_i$ bijectif ou si on veut $\mathrm{gr}(u)$ bijectif, ce qui est
immédiat et sauf erreur figure déjà en lemme quelque part au Chap III
(j'espère, avec la généralité qui convient, et bien entendu modulé \ldots).
\add{cf diagramme}

\begin{tikzcd}
\mathrm{gr}\,B \arrow[r, "\mathrm{surj}"] & \mathrm{gr}\,C \\
B_0 \otimes \mathrm{gr}\,A_i \arrow[r, "\mathrm{bij}"] \arrow[u, "\mathrm{surj}"] & C_0 \otimes \mathrm{gr}\,A_i \arrow[u, "\mathrm{bij}"]
\end{tikzcd}

\note{le diagramme est de sa main, tracé au bas du feuillet sous la dernière
ligne dactylographiée.}

\page{10}\note{sa page 25}

Enfin, il nous sera commode d'avoir un critère permettant de déduire de la
propriété P une propriété de platitude. Nous utiliserons :

\textbf{Lemme 3.3.} Soient $A$ un anneau, $(P_i)_{i \in I}$ une famille de
$A$-modules, \add{$P$ leur produit} $M$ un $A$-module de présentation finie.
Alors l'homomorphisme canonique
\[
  M \otimes P \longrightarrow \prod_i M \otimes P_i
\]
est un isomorphisme.

En effet, écrivant $M$ comme \struck{comme} un conoyau d'un homomorphisme de
modules libres \add{de type fini} $L' \to L$, et notant que le foncteur
produit est exact donc exact à gauche, on est ramené au cas où $M = A^n$, $n$
entier $> 0$, où c'est trivial.

\textbf{Corollaire 3.4.} Supposons $A$ tel que tout module de type fini sur
$A$ soit de présentation finie, par exemple $A$ nœthérien. Alors tout produit
de $A$-modules plats est un $A$-module plat.

En effet, il suffit de montrer que $P \otimes$ transforme un monomorphisme
$M \to N$ en un monomorphisme, quand $M$ et $N$ sont de type fini ; le lemme 3
donne alors immédiatement le résultat voulu.

\textbf{Corollaire 3.5.} Supposons que $A$ satisfasse \struck{les conditions}
\add{par exemple $A$ nœthérien} de 3.4. Soit $P$ un $A$-module
\struck{topologique} isomorphe à \struck{une} la limite projective
\struck{\ill{}} d'un système projectif strict, \struck{de $A$} admettant un
système cofinal dénombrable, de $A$-modules projectifs. Alors $P$ est un
$A$-module plat.
\note{le passage biffé au milieu de l'énoncé est frappé de x et couvre deux
lignes et demie ; on y lit encore « de $A$-modules projectifs de type fini
(``stricte'' = par un système projectif strict\ldots\ ``dénombrable'' par un
système proj\ldots) ». Les mots « stricte » et « dénombrable » sont repris à la
main au-dessus.}

En effet, on voit facilement que sous les conditions dites, $P$ est isomorphe
à un produit de $A$-modules projectifs, donc plats, d'où le résultat.

\page{11}\note{sa page 26}

\textbf{Théorème 3.6.} Soient $A$ un anneau préadmissible muni d'un idéal de
définition $I = I_0$ et d'un système fondamental $(I_{i})$ d'idéaux de
définition $\subset I_0$, $B$ une $A$-algèbre topologique. On suppose les
$A_i = A/I_iA$ nœthériens, et de plus les conditions suivantes sur $B$ :

\begin{enumerate}
  \item[a)] \struck{les $A_i$ sont nœthériens}
  \item[a)] Les \struck{Module} Modules topologiques \add{$B$ sur $A$}
    \struck{\ill{}} satisfont la condition P.
  \item[b)] \add{$B_0$ est formellement simple sur $A_0$.}
  \item[c)] $B$ est adique \struck{pour un idéal contenant $IB$}
  \item[d)] Les idéaux $I_iB$ dans $B$ sont fermés (\struck{inutile}
    \add{automatiquement vérifié} si $B$ est nœthérien).
\end{enumerate}
\note{la liste est renumérotée à la main : l'ancien a) est biffé, une clause
b) entièrement manuscrite est intercalée, et l'ancien b) devient c).}

Sous ces conditions, on a ce qui suit :

\begin{enumerate}
  \item[(i)] $B$ est formellement simple sur $A$.
  \item[(ii)] Soit $C$ une $A$-algèbre topologique satisfaisant à c) et d)
    (mais pas nécessairement à a)), et soit $u_0 : B_0 \to C_0$ un
    isomorphisme d'algèbres topologiques. Alors il existe un
    $A$-homomorphisme continu $u : B \to C$ se réduisant suivant $u_0$ (en
    vertu de (i) et de 1.7.). Ce dernier étant choisi, les conditions
    suivantes sont équivalentes :
\end{enumerate}

\begin{enumerate}
  \item[1)] $B$ est formellement simple sur $A$.
  \item[2)] \struck{\ill{}} $B$ satisfait la condition a) ci-dessus, i.e.
    \struck{\ill{}} satisfait la condition P.
  \item[3)] Les $B_i$ sont plats sur $A_i$.
  \item[4)] L'homomorphisme $u : B \to C$ est un isomorphisme de
    $A$-algèbres top.
  \item[5)] \add{$C$ est isom.\ à $B$ comme $A$-algèbre top.}
\end{enumerate}

Démonstration. L'assertion (i) résulte déjà de a), grâce à 3.1. Prouvons donc
(ii). En vertu de (i) déjà prouvé, on a 2) $\Rightarrow$ 1), et évidemment
4) $\Rightarrow$ 2), d'ailleurs 2) $\Rightarrow$ 3) en vertu de 3.5., grâce
au fait que les $A_i$ sont nœthériens. Reste à prouver 1) $\Rightarrow$ 4) et
3) $\Rightarrow$ 4).

Supposons 1) vérifié, alors en vertu de 1.7. $u_0^{-1} : C_0 \to B_0$ se
prolonge en un homomorphisme d'algèbres top $v : C \to B$ (grâce aux
hypothèses c) et d) ). Or on a déjà noté que les $B_i$ sont plats sur les
$A_i$. Le lemme 3.2. (i) et (iii) (où on échange les rôles de $B$, $C$)

\page{12}\note{sa page 27}

garantit alors que $v$ est un isomorphisme d'algèbres topologiques (grâce à c)
et d) pour $B$ et $C$). Donc \struck{les} on a 2), \add{donc on a 3)}
\struck{et} le même argument en sens inverse montre que \struck{\ill{}}
\add{3) implique 4). --- Enfin 4) $\Rightarrow$ 5) et 5) $\Rightarrow$ 2).}
On peut certainement améliorer la rédaction de la démonstration, évitant
certaines redites.

Pour simplifier, nous nous restreindrons par la suite à des hypothèses
nœthériennes (permettant de nous débarrasser en particulier de la désagréable
hypothèse d) sur $B$).

\textbf{Définition 3.7.} Soit $A \to B$ un homomorphisme continu d'anneaux
nœthériens préadiques. On dit que $B$ est une \emph{algèbre (topologique) de
Cohen} sur $A$ si \add{a)} $B$ est séparé et complet, i.e.\ \add{un} un
anneau adique, \struck{\ill{}} \add{b)} s'il satisfait à la condition P, et
si enfin \add{c)} $B$ est formellement \struck{\ill{}} \add{simple} sur $A$,
\add{ou encore, $I$ étant un idéal de définition de $A$, $B_0 = B/ID$
\ill{} formellement simple sur $A_0$}.
\note{la phrase manuscrite qui ferme la définition court sur deux lignes en
interligne ; elle porte en tête un « $= A/I$ » et sa fin n'est lisible qu'en
partie.}

On notera qu'en vertu de 3.1., si $I$ est un idéal de définition de $A$, la
dernière condition équivaut à la condition que $B/IB$ est formellement simple
sur $A/IA$. D'ailleurs, on peut interpréter 3.6. de la façon suivante (et il
est \struck{\ill{}} plus raisonnable de reléguer 3.6. en lemme, et de réserver
le corollaire en théorème) :

\textbf{Corollaire 3.8.} Soient $A$ un anneau \struck{préadique} nœthérien,
$I$ un idéal de $A$, munissons $A$ de la topologie $I$-préadique, posons
$A_0 = A/IA$, et soit $B_0$ une $A_0$-algèbre topologique. \struck{\ill{}}
\add{Supposons qu'}il existe une algèbre de Cohen $B$ sur $A$ telle que $B/IB$
soit isomorphe à $B_0$ comme $A_0$-algèbre topologique, \add{ce qui implique
$B_0$ \uncertain{aussi de Cohen} sur $A_0$}, alors $B$ est unique à \add{un}
isomorphisme près \struck{induisant l'isomorphisme donné sur} (compatible avec
les isomorphismes donnés des algèbres \struck{réduites} réduites mod $I$ avec
$B_0$). De plus, si $B$ est une $A$-algèbre topologique adique telle que
$B/IB$ soit isomorphisme à $B_0$ comme $A_0$-algèbre topologique, alors les
conditions suivantes sur $B$ sont équivalentes : 0) $B$ est une $A$-algèbre de
Cohen
\note{« soit isomorphisme à $B_0$ » est frappé ainsi ; la lettre qui ouvre la
liste des conditions équivalentes est repassée à l'encre et se lit 0) ou a).}

\page{13}\note{sa page 28}

\begin{enumerate}
  \item[1)] $B$ est formellement simple sur $A$
  \item[2)] $B$ satisfait la condition P
  \item[3)] $B$ est plat sur $A$
  \item[3 bis)] Pour tout entier $n \geqslant 0$, $B_n = B/I^{n+1}B$ est plat
    sur $A_n = A/I^{n+1}$. \struck{\ill{}}
\end{enumerate}

Il reste seulement à montrer l'équivalence de 3) et 3 bis), qui est contenue
en effet (moyennant les hypothèses nœthériennes faites) dans la platitude
supérieure.

\struck{Le résultat précédent} Il conviendrait de noter \struck{aussi}
\add{dès après 3.7.} que ``algèbre de Cohen'' est une notion stable par
extension de la base via produits tensoriels complétés : c'est une trivialité.

Le résultat 3.8. montre l'intérêt qu'il y a, $A$, $I$ et $B_0$ étant donnés,
de disposer d'un \emph{théorème d'existence} pour une algèbre de Cohen $B$ sur
$A$ se réduisant suivant $B_0$. On notera que la construction habituelle de
proche en proche nous ramène \struck{au cas où $A$ est discret} (remplaçant
$A$ par des $A_i$ et $I$ par des puissances) au cas où $I$ est de carré nul,
(donc où $A$ est discret). On obtient alors un problème de nature
essentiellement infinitésimale, dont nous ignorons s'il existe toujours une
solution, faute de savoir caractériser la structure des algèbres de Cohen
$B_0$ sur l'anneau discret $A_0$. Cependant, nous obtenons \add{(3.10.)} un
théorème d'existence dans un cas particulier intéressant, en utilisant les
résultats du Nº2. Notons d'abord

\struck{\ill{}}
\note{un long passage biffé de deux grands traits diagonaux, sur six lignes :
il portait un « Théorème 3.9 » sur $A$ local nœthérien de corps résiduel $k$,
$B_0$ une $k$-algèbre qui soit un anneau local nœthérien complet de corps
résiduel $K$ de degré de multiplicité radicielle finie sur $k$, extension
séparable de $k$, et s'achevait sur « $B_0$ est (pour sa topolo- ». Il est
remplacé par le corollaire qui suit.}

\textbf{Corollaire 3.9.} Soit $A \to B$ un homomorphisme local d'anneaux
locaux nœthériens, \struck{tel que l'extension résiduelle $k(B)/k(A)$ soit à
multiplicité}

\page{14}\note{sa page 29}

radicielle finie (par exemple, \struck{soit} séparable) \add{Conditions
équivalentes :} \add{1)} Pour que $B$ \struck{soit} \add{est} une algèbre de
Cohen sur $A$ (sous-entendu : pour les topologies habituelles, --- il faudrait
le préciser dans 3.7.), \struck{il f et suffit que} $B$ \add{2)}
\struck{est} plat sur $A$, complet, et \struck{que} $B_0 = B/\mathfrak{m}B$
\struck{soit} \add{est} une algèbre de Cohen sur $k = A/\mathfrak{m}$, i.e.
formellement simple sur $k$, \add{3) $B$ est formellement \uncertain{simple}
sur $A$} \struck{\ill{}}
\note{l'énoncé est retourné à la main en liste de conditions équivalentes :
trois numéros manuscrits 1), 2), 3) sont plantés dans la phrase, la troisième
condition est écrite en entier à l'encre puis en partie biffée d'un
griffonnage, et « il faut et suffit que » est biffé.}

(On aura noté \struck{\ill{}} après 3.7. que pour une algèbre topologique sur
un corps $k$, Cohen $=$ adique et formellement simple).
\struck{La nécessité} étant déjà \struck{acquise} (sans hypothèse de
multiplicité radicielle finie), prouvons \struck{la suffisance}
\add{1) $\Rightarrow$ 2)} \add{2) $\Rightarrow$ 1)}. En vertu de la
définition, on est ramené simplement à prouver que $B$ satisfait la condition
P. Cela nous ramène au cas où $A$ est artinien. En vertu du Nº2, on sait que
$B_0$ est régulier, considérons un système régulier de paramètres de $B_0$
(un système de paramètres formant une $B_0$-suite suffirait), relevons le en
un système de paramètres de $B$, soit $x = (x_1, \ldots, x_r)$. Alors les
idéaux engendrés par les $x^n = (x_1^n, \ldots, x_r^n)$, pour $n \geqslant 0$,
forment un système fondamental de voisinages \add{$J_n$} de $0$ dans $B$,
d'autre part on constate aussitôt que les $B/J_n$ sont des modules plats sur
$A$ (il y a un \struck{lemme} corollaire dans la platitude supérieure, dans le
cas où on divise par un seul élément \struck{\ill{}} régulier, cas auquel on
est ramené aussitôt). Comme $A$ est artinien, ils sont donc projectifs.

\add{\ill{}}
\note{une phrase manuscrite traverse ici la page sur deux lignes, entre
« projectifs. » et le NB, et déborde des deux côtés ; on n'en lit avec
certitude que des fragments — un « que \ldots\ 3) $\Rightarrow$ 2) », un
« \ldots\ suffisante », et la fin « \ldots\ d'ailleurs ci-dessous 3.11. » La
ligne dactylographiée « projectifs. » porte à sa suite un mot biffé.}

NB On n'a utilisé l'hypothèse sur la multiplicité radicielle que pour
\struck{pouvoir} savoir que $B_0$ est Cohen-Macaulay. L'hypothèse faite est
très probablement inutile. Serre vient de me dire que Mac-Lane a prouvé la
réciproque de 2.3. sans hypothèse de multiplicité rad finie, dans ``Note on a
relative structure of $p$-adic fields'' Ann Math t.41 1940, page 751--753.
Cela vaudrait peut-être le coup, pour \struck{liquider} le coup d'œil, que
tu regardes cela et donnes le résultat complet dans 2.3. On
\note{tout le NB est barré de six longs traits diagonaux au crayon bleu et
enfermé dans un crochet ; il est barré, non supprimé, et reste lisible d'un
bout à l'autre.}

\page{15}\note{sa page 30}

\struck{pourra alors resper} \add{fait donc dans 3.9.} l'hypothèse
\struck{\ill{}} avec : ``ou \struck{\ill{}} $r(B) = r(A)B$''
\struck{\ill{}}

\textbf{Corollaire 3.10.} Soient $A$ un anneau local nœthérien, de corps
résiduel $k$, $B_0$ une $k$-algèbre qui soit un anneau local nœthérien, et
soit une $k$-algèbre de Cohen i.e.\ une $k$-algèbre formellement simple
\struck{pour sa} et complète pour sa topologie ordinaire. Supposons de plus
l'extension résiduelle $K/k$ séparable (ce qui n'est donc pas une restriction
si $B_0 = K$ \struck{\ill{}} mais en est évidemment une en général). Alors il
existe une $A$-algèbre topologique de Cohen \add{$B$} qui se réduit suivant
$B_0$, $B$ est unique à isomorphisme près en vertu de 3.8., et c'est un anneau
nœthérien local complet comme il se doit.

En vertu du théorème de structure 3.4. on a $B_0 \simeq K[[T_1, \ldots, T_n]]$
et on est ramené, grâce au critère 3.9. 2) et la transitivité de la platitude,
aux deux cas types : $1^{\mathrm{o}}$) $B_0 = K$, $2^{\mathrm{o}}$)
$B_0 = k[[T_1, \ldots, T_n]]$. Dans le premier cas, on utilise le résultat
général de $0_{\mathrm{III}}$ \add{10.3.1.}, et dans le deuxième cas, il
suffit de prendre $B = \hat{A}[[T_1, \ldots, T_n]]$.

\struck{Théorème 3.11} \textbf{Corollaire 3.11.} Sous les conditions de 3.9.,
\add{si $k(B)/k(A)$ de mult.\ rad.\ finie,} les conditions 1) et 2)
équivalent aussi à : 3) $B$ est complet et formellement simple sur $A$.

En effet, on sait déjà que \struck{c'est nécessaire} reste à prouver
\struck{la} 3) $\Rightarrow$ 1) \add{1) $\Rightarrow$ 3)}
\struck{suffisance} ou encore 3) $\Rightarrow$ 2), i.e.\ on est ramené à
prouver simplement que $B$ est plat sur $A$. Or \struck{\ill{}} faisant une
extension de la base $A \to A'$, $A'$ local \add{de type fini} libre sur $A$ à
extension radicielle convenable (cf $0_{\mathrm{III}}$), on trouve un $B'$
\add{local} sur $A'$ qui sera encore formellement simple sur $A'$, et de plus
à extension résiduelle \struck{triviale} séparable. On dispose alors du
théorème d'existence 3.10. qui nous garantit, grâce à 3.8. (équiv de 0) et
1) ) que $B'$ est Cohen sur $A'$ donc plat sur $A'$. Donc $B$ est plat sur
$A$.

\marginal{Réduction inutile dans le cas ``vrai'' du NB}
\note{en bas à droite, de sa main : elle renvoie au NB barré de la page
précédente.}

\page{16}\note{sa page 31}

En particulier, on obtient :

\textbf{Corollaire 3.12.} Soient $A \to B$ un homomorphisme local d'anneaux
locaux nœthériens, tels que $r(B) = r(A)B$ et \struck{que l'extens} soit
$K/k$ l'extension résiduelle. Conditions équivalentes : 1) $B$ est complet et
formellement simple sur $A$. 2) $B$ est complet, plat sur $A$, et $K/k$ est
séparable. 3) $B$ est une algèbre de Cohen sur $A$. De plus, une extension
séparable $K$ de $k$ étant donnée, il existe une algèbre de Cohen $B$ sur $A$
se réduisant suivant $K$, et $B$ est unique à isomorphisme de $A$-algèbres
près.

La rédaction commence à faire un peu touffu, j'espère que tu pourras la rendre
plus lisible. D'ailleurs, je me repens un peu de la définition 3.7. telle
quelle, il n'est pas vraiment naturel de supposer $B$ séparée et complète pour
sa topologie propre, mais simplement pour sa topologie $IB$-adique (ce qui est
moins fort) ; mais une telle modification demanderait des réajustements
analogues pour 3.2., 3.6. et la suite. Je t'en laisse donc le choix, pensant
que le choix adopté n'aura pas de grande importance pratique pour la suite :
dans 3.9. à 3.12. cela obligera simplement à supprimer à certains endroits le
mot ``complet'', \add{ou à le remplacer par un autre complet}
\struck{en particulier ``Cohen'' deviendra synonyme dans le contexte des
anneaux locaux nœthériens de ``formellement simple''}, --- sous réserve de
l'embêtante hypothèse de \struck{multiplicité radicielle}.
\struck{Il est vrai} Ce serait plutôt un avantage, \struck{mais il est vrai
par contre que cela ne s'harmoniserait plus} de plus les \struck{anneaux de
Cohen sont}

\add{\ill{}}
\note{le bas du feuillet porte, sur quatre lignes et dans une écriture très
rapide, un début de « Corollaire 3.13 » manuscrit — on y lit « $A \to B$
\ldots\ local d'anneaux locaux nœth., \ldots\ de corps \ldots\ à mult.\ deg.\
finie. Alors conditions équivalentes » — aussitôt annulé par cinq grandes
boucles. C'est le premier jet du corollaire que le feuillet suivant reprend
au propre.}

\page{17}\note{sa page 32}

\textbf{Corollaire 3.13.} Soit $A \to B$ un homomorphisme local d'anneaux
locaux nœthériens, on suppose $k(B)/k(A)$ à multiplicité radicielle finie
(par exemple une extension séparable, ou de type fini), \add{ou
$k(B) = k(A)B$.} \struck{Conditions} On munit $A$ et $B$ de leurs topologies
habituelles. Alors les conditions suivantes sont équivalentes : (i) $B$ est
formellement simple sur $A$ (ii) $B$ est plat sur $A$, et
$B_0 = B \otimes_A k$ est une algèbre universellement régulière sur $k$
(cf 2.6.).

En effet, les conditions (i) et (ii) étant invariantes par changement de $B$
en $\hat{B}$, on est ramené au cas où $B$ est complet, où cela est contenu
dans 3.11. \add{resp.\ 2.12.} compte tenu que pour $B_0$, formellement simple
$=$ universellement régulier en vertu de \uncertain{2.6.}

Du point de vue pratique, les résultats du présent numéro consistent
essentiellement en 3.13., \struck{et l'existence} l'unicité \struck{\ill{}}
(moyennant les conditions de 3.13.) du $B$ correspondant à un $B_0$ donné,
sous réserve que $B$ soit complet pour la topologie $r(A)$-adique (cette
\struck{\ill{}} unicité étant un corollaire immédiat d'ailleurs, puisqu'on
dispose pour $B$ des deux propriétés complémentaires formellement simple $+$
plat), \add{et enfin l'existence moyennant la séparabilité de $K/k$.} À force
de dévissage et de généralité, cette substance n'apparait pas assez
clairement, et la terminologie ``algèbre de Cohen'' introduite concurremment à
``formellement simple'' risque \struck{\ill{}} \add{plutôt} d'embrouiller les
idées (d'autant plus que nous n'utiliserons plus jamais que ``formellement
simple''). Aussi le présent numéro rendra peut-être plus service si on
commence à énoncer simplement 3.13. en théorème, avec le corollaire bien senti
précédent, et qu'on les démontre simultanément en deux pages ou trois,
\struck{avec} en laissant de côté tout le fatras anneaux admissibles, algèbres
de Cohen etc. Commencer par l'existence d'un $B$ satisfaisant à la fois (i) et
(ii) si $K/k$ est séparable, se ramener à ce cas pour 3.13. par descente comme
dans 3.11. prouver 3.13. alors par comparaison du $B$-type avec le $B$ donné,
ce qui donne en même temps l'unicité, et un point c'est tout.

\marginal{(garder cependant 3.1. au Nº3. On peut signaler 2.6. \ill{} pour
remarque, à réserver pour l'avenir)}
\note{écrit verticalement dans la marge de gauche, sur quatre lignes, en
regard du dernier alinéa.}

\section{Questions ouvertes}

\note{Le feuillet 18 est une couverture : une feuille par ailleurs nue,
portant en haut à droite, de sa main, le seul mot « Questions ». Il ne reçoit
pas de \texttt{page} ; les deux feuillets manuscrits qu'il
enveloppe suivent. Ceux-ci sont d'une écriture rapide et fine, et la prose de
liaison ne se laisse lire qu'en partie : ce qui suit garde les énoncés et
laisse \ill{} partout où la feuille ne porte pas de lecture sûre.}

\page{19}

Questions ouvertes

\begin{enumerate}
  \item[1)] Soit $A$ algèbre locale nœthérienne sur $k$. Est-il vrai
    que $A$ formellement simple sur $k$ si $A$ est universellement régulière
    sur $k$ [\uncertain{définition}, pour hom.\ local $A \to B$ d'anneaux
    locaux nœthériens, équivalence entre « ft simple »
    et « universellement régulier » \ill{} sur $A$ ].
  \item[2)] Soit $A$ algèbre locale nœthérienne sur $k$, ft
    simple sur $k$ (pour sa topologie habituelle).
    Est-il vrai que $A$ est ft simple pour sa \uncertain{top.}
    \ill{} discrète ? \uncertain{Ou} : $A$ est complète ? (\ill{}
    possible que le \uncertain{rigoureux} soit \ill{} \ill{}
    \ill{} \ill{} $0$). --- On peut \ill{}
    \ill{} \ill{}. $A$ plat sur $\Lambda$,
    \ill{} $\Lambda$ \ill{} nœthériens : si les \ill{}
    de $A$ est ft simple sur \ill{} de $\Lambda$
    \ill{} pour une topologie habituelle, $A$ est-il ft.\
    simple sur $\Lambda$ pour sa top.\ discrète.
  \item[3)] $B$ plat sur $A$, $A$ et $B$ nœthériens.
    L'\ill{} \ill{} de $\mathrm{Spec}$ \uncertain{$B$} dans l'\ill{}
    \ill{} fermé dans le $\mathrm{Spec}\,A$ \ill{} est ft simple
    sur le corps résiduel de $A$, est-il ouvert ?
    \struck{\ill{}} [ Si $B$ \ill{} pas \ill{} $A$,
\end{enumerate}

\marginal{\ill{}}
\note{un long bloc de sa main court en oblique, de bas en haut, dans toute la
marge de gauche du feuillet, sur une douzaine de lignes ; c'est un ajout
postérieur au corps de la page. Il n'a pas pu être lu. On y distingue
seulement, dans les premières lignes, « Faut-il supposer \ill{} vraiment
\ill{}\ ? Peut-être qu'\ldots », ainsi qu'une formule « $\Lambda = VA$ » ou
« $\mathfrak{p} \cap A$ » ; à hauteur de la question 1) figure en outre une
brève note isolée, de deux ou trois caractères.}

\page{20}

l'\ill{} des Vs de tout $B$ \ill{} $B$ l'\uncertain{est} aussi ? ].
Il \uncertain{faut-il} \uncertain{supposer} \ill{} que $A$ et $B$
soient des « bons anneaux ».

\struck{l'\ill{} $A$, $A$ \ill{} local, $A \to B$ localisation ?}
\marginal{Cf IV 2.10.1}
\note{la ligne barrée et la mention « Cf IV 2.10.1 » sont réunies par un
crochet tracé à l'encre ; c'est le seul renvoi précis à un autre chapitre que
portent ces deux feuillets.}

\ill{} d'extension résiduelle finie.

\ill{} \uncertain{contre-t-il} critère jacobien de \ill{} de
Nagata pour une \uncertain{algèbre} locale nœth \ill{} \ill{} on a dit
\ill{} en \uncertain{vrais}.

\ill{} Ok.\ en caractéristique $0$ pour les \textbf{bons}
\ill{} \ill{}, \uncertain{puis} : \textbf{la résolution des
singularités}.
\note{« bons », « résolution des singularités » et le dernier mot sont
soulignés sur la feuille.}

\begin{enumerate}
  \item[4)] $A \to B$ hom d'anneaux nœth. L'\ill{} des Vs de
    tout $B$ \ill{} le quotient $D$ \ill{} sur $A$ \ill{}
    est-il ouvert ? [ \ill{} \ill{} enfin
    que \ill{} $VA$ \uncertain{$\mathfrak{p} \cap \mathrm{Spec}\,A$},
    et \ill{} $VA$ \ill{} de \ill{} $VA$ \ill{}
    \ill{} \uncertain{$A' = A_{\mathfrak{p}}$}, \ill{}
    \uncertain{$B' = A_{\mathfrak{p}}$} \ill{} \uncertain{$\Lambda \subset B'$}
    \ill{} \ill{} $B \otimes_{A'}$ \ill{} \ill{}
    \ill{} \ill{} $B'$ \ill{} plat sur $A'$. ]
\end{enumerate}
\note{la question 4) est tracée dans la moitié droite du feuillet, en colonne
étroite, et ses dernières lignes chevauchent l'en-tête imprimé du papier ;
elles ne se lisent que par fragments. Les identifications $A'$, $B'$ portées
ici sont douteuses.}

\marginal{\ill{}}
\note{comme au feuillet précédent, un bloc écrit en oblique occupe toute la
marge de gauche, plus dense encore : une vingtaine de lignes, avec au moins
trois renvois numérotés (1), (2), (3) et des passages soulignés. Il n'a pas pu
être lu.}

\end{document}
